\section{重要矢量恒等式与数学工具}
%=============================================================================

推导过程中使用的主要矢量恒等式汇总如下，方便自查：

\begin{align}
  \nabla\times(\psi\mathbf{A}) &= \psi(\nabla\times\mathbf{A}) + \nabla\psi\times\mathbf{A} \label{eq:vec1}\\
  \nabla\times\nabla\times\mathbf{A} &= \nabla(\nabla\cdot\mathbf{A}) - \nabla^2\mathbf{A} \label{eq:vec2}\\
  \nabla\times(\mathbf{r}\psi) &= \nabla\psi\times\mathbf{r} = -\mathbf{r}\times\nabla\psi \quad (\text{因 }\nabla\times\mathbf{r}=0) \label{eq:vec3}\\
  \mathbf{r}\cdot\nabla^2\mathbf{E} &= \nabla^2(\mathbf{r}\cdot\mathbf{E}) - 2\nabla\cdot\mathbf{E} \label{eq:vec4}\\
  \int \mathbf{A}\cdot(\nabla\times\mathbf{B})\,d^3r &= \int (\nabla\times\mathbf{A})\cdot\mathbf{B}\,d^3r
    \quad (\text{边界项为零}) \label{eq:vec5}
\end{align}

以及球谐函数的梯度关系：
\begin{align}
  \nabla Y_{lm} &= \frac{1}{r}\Bigl(\hat{\boldsymbol{\theta}}\,\partial_\theta Y_{lm}
    + \hat{\boldsymbol{\phi}}\,\frac{1}{\sin\theta}\,\partial_\phi Y_{lm}\Bigr) \label{eq:gradY}
\end{align}

和矢量球谐函数 $\mathbf{X}_{lm}$ 的分量形式：
$$
\mathbf{X}_{lm}(\theta,\phi) = \frac{1}{\sqrt{l(l+1)}}\Bigl(
  i\hat{\boldsymbol{\theta}}\,\frac{m}{\sin\theta}P_l^m(\cos\theta)
  - \hat{\boldsymbol{\phi}}\,\frac{d}{d\theta}P_l^m(\cos\theta)
\Bigr)e^{im\phi}
$$

球 Bessel 函数满足的微分方程：
$$
\Bigl[\frac{d^2}{dr^2} + \frac{2}{r}\frac{d}{dr} + \Bigl(k^2 - \frac{l(l+1)}{r^2}\Bigr)\Bigr] j_l(kr) = 0
$$

由此可推导 Riccati--Bessel 函数 $\mathfrak{g}_l(kr)=kr\,j_l(kr)$ 满足的方程：
$$
\Bigl(\frac{d^2}{dr^2} + k^2\Bigr)\mathfrak{g}_l(kr) = \frac{l(l+1)}{r^2}\,\mathfrak{g}_l(kr)
$$

\vspace{1em}
\begin{center}
  \fcolorbox{annotblue!50}{annotblue!5}{\parbox{0.85\textwidth}{\centering
  \textcolor{annotblue}{$\blacktriangleright$} \textbf{结论}：Eqs.(15)--(16) 是 Grahn 2012 框架的\textbf{核心实用公式}，\\
  将 $a_E/a_M$ 的计算转化为散射体体积内的 $\mathbf{J}_S$ 积分，\\
  避免了球面积分和数值求导，适合阵列中的逐个粒子分析。}}
\end{center}

%=============================================================================
